\documentclass[11pt]{article}
\usepackage[a4paper,margin=20mm]{geometry}
\usepackage{amsmath,amssymb,amsthm,microtype}
\usepackage[colorlinks=true,linkcolor=blue,citecolor=blue,urlcolor=blue]{hyperref}
\hypersetup{pdftitle={Finite gate alphabets, scalar phase, and auxiliary erasure},
  pdfauthor={TNLean contributors}}
\newtheorem{theorem}{Theorem}
\newtheorem{lemma}[theorem]{Lemma}
\newtheorem{corollary}[theorem]{Corollary}
\title{Finite gate alphabets, scalar phase, and auxiliary erasure}
\author{TNLean contributors}
\date{3 October 2026}
\begin{document}
\maketitle

The exact circuit theorem for bounded-cut-rank unitaries permits arbitrary
one- and two-qubit gates, with every auxiliary initialized and returned to
zero; see \cite[Section~5]{exact-mpu-circuits}. The present note
examines conversion to a fixed finite gate alphabet while retaining the
actual scalar phase of the physical unitary. A finite gate alphabet is a
restriction on circuit gates, distinct from a finite set of MPU coefficient
values.

\textbf{Mathematical status.} The arguments below are mathematical proofs.
The finite-alphabet conversion and the exact-erasure results are not
asserted to be formalized in Lean. The arbitrary-gate circuit theorem used
in the MPU application is proved in the cited note and formalized in
\nolinkurl{TNLean/MPS/MPU/ExactBoundedCutRankCircuit.lean}.
All norms below are Euclidean operator norms, including the norm of a
rectangular map from physical inputs to the complete register.

\section{Approximation with scalar phase retained}

Let $m\geq N\geq1$, let $U$ act on $N$ qubits, and let $J:\mathbb C^{2^N}\to\mathbb C^{2^m}$
initialize $m-N$ auxiliary qubits in their computational zero state, and
suppose that an exact circuit $C=V_G\cdots V_1$ of $G\geq1$ arbitrary
one- and two-qubit gates satisfies
\begin{equation}\label{eq:finite-gate-exact-input}
 CJ=JU.
\end{equation}
The finite alphabet will be
\begin{equation}\label{eq:finite-gate-alphabet}
 \mathcal A=\{H,T,T^{-1},\operatorname{CNOT}\},\qquad
 H=\frac1{\sqrt2}\begin{pmatrix}1&1\\1&-1\end{pmatrix},\quad
 T=\operatorname{diag}(1,e^{i\pi/4}).
\end{equation}
The inclusion of $T^{-1}$ changes gate counts by at most a constant, since
$T^{-1}=T^7$. Placements of these gates on the register are allowed as usual.

\begin{theorem}[Finite-alphabet approximation with retained scalar phase]
\label{thm:finite-gate-phase-lift}
For $0<\varepsilon<1$, there is an $\mathcal A$-circuit $C_\varepsilon$
on $m+1$ qubits, with arbitrary pairs allowed to interact, such that
\begin{equation}\label{eq:finite-gate-input-error}
 \|C_\varepsilon J_+-J_+U\|\leq\varepsilon,
 \qquad J_+\psi=|0\rangle\otimes J\psi.
\end{equation}
Its gate count is
\begin{equation}\label{eq:finite-gate-count}
 O\left(G\bigl[1+\log(G/\varepsilon)\bigr]^c\right),
 \qquad c=\frac{\log5}{\log(3/2)}.
\end{equation}
The constant is absolute. In particular, the scalar phase of $U$ is retained.
The norm of the component outside the fully initialized auxiliary subspace
is at most $\varepsilon$; its probability is at most $\varepsilon^2$ on
every unit physical input. Exact auxiliary erasure is not asserted.

If the original gates join neighboring qubits on a line, the same gate bound
holds with neighboring gates from $\mathcal A$ after adding $m$ phase
qubits initialized to zero. Thus the auxiliary count increases from $m-N$ to $2m-N$.
Alternatively, one additional phase qubit initialized to zero suffices on
a line with gate count
$O\bigl(G[m+(1+\log(G/\varepsilon))^c]\bigr)$.
\end{theorem}

The proof establishes density in special unitary groups in operator norm,
including the scalar phase. The needed quantitative approximation statement is
proved next, to specify this distinction and the precision dependence.

\begin{lemma}[Special-unitary words from the stated alphabet]
\label{lem:finite-gate-su-density}
For $k=2,3$, there is a finite inverse-closed set $\mathcal S_k\subset
\operatorname{SU}(2^k)$, each member an exact constant-length
$\mathcal A$-word on $k$ qubits, whose generated subgroup is dense in
$\operatorname{SU}(2^k)$ in operator norm.
\end{lemma}

\begin{proof}
Put $\zeta=e^{i\pi/4}$. Direct multiplication gives
\begin{equation}
 (HT^2)^3=\zeta I_2.
\end{equation}
Consequently the following are exact alphabet words, including their phases:
\begin{equation}
 P=\zeta^{-1}THTH,\qquad Q=HPH.
\end{equation}
Both have determinant one. With $a=\pi/8$, $s=\sin a$, and $t=\cos a$,
one has
\begin{equation}
 P=t^2I-i\bigl(stX+s^2Y+stZ\bigr),\qquad
 Q=t^2I-i\bigl(stX-s^2Y+stZ\bigr).
\end{equation}
Their rotation axes are not parallel. If the eigenvalues of $P$ were roots
of unity, their sum $\operatorname{Tr}P=1+1/\sqrt2$ would be an algebraic
integer. Its norm from $\mathbb Q(\sqrt2)$ is $1/2$, so it is not an
algebraic integer. Thus the powers of $P$ are dense in its one-parameter
rotation group; the same holds for $Q$. The two infinitesimal axes and their
commutator span $\mathfrak{su}(2)$. The product and commutator limits for
matrix exponentials then show that the closed group generated by $P,Q$
contains every exponential in $\mathfrak{su}(2)$, hence all of
$\operatorname{SU}(2)$.

On $k$ qubits take the copies $P_i,Q_i$, their inverses, and the words
\[
 \operatorname{CNOT}_{ij}P_j\operatorname{CNOT}_{ij},\qquad
 \operatorname{CNOT}_{ij}Q_j\operatorname{CNOT}_{ij},
\]
together with their inverses, for $i\ne j$. All these words have
determinant one, also when an individual CNOT does not. Their closure contains every local special unitary and, since
\begin{equation}
 \operatorname{CNOT}_{ij} Z_j\operatorname{CNOT}_{ij}=Z_iZ_j,
\end{equation}
every one-parameter rotation generated by $Z_iZ_j$. Local conjugations give
all two-site Pauli products. For $k=3$, commutators of products supported on
$\{i,j\}$ and $\{j,h\}$, with different Pauli factors at $j$, give all
three-site Pauli products. Together with the local Pauli operators these
form a basis of the traceless Hermitian matrices. The product and commutator
limits therefore generate every exponential of a traceless skew-Hermitian
matrix. Every special unitary has such a logarithm, by choosing eigenangles
with sum zero. This proves the claim for $k=2,3$.
\end{proof}

\begin{lemma}[Quantitative approximation in fixed dimension]
\label{lem:finite-gate-quantitative-su}
Let $q$ be fixed, and let $\mathcal S\subset\operatorname{SU}(q)$ be a
finite inverse-closed set generating a dense subgroup. Every
$W\in\operatorname{SU}(q)$ admits an $\mathcal S$-word $W_\delta$ such
that $\|W_\delta-W\|\leq\delta$, with word length
$O_{q,\mathcal S}\bigl((1+\log(1/\delta))^c\bigr)$ for $0<\delta<1$,
where $c=\log5/\log(3/2)$.
\end{lemma}

\begin{proof}
The balanced commutator construction is recorded first. If
$\Delta\in\operatorname{SU}(q)$ is sufficiently close to $I$, its
principal logarithm is $\Delta=e^{iH}$ with $H$ traceless Hermitian and
$\|H\|\leq2\|\Delta-I\|$. Tracelessness follows because the eigenangles
are small and their sum belongs to $2\pi\mathbb Z$.
There is an orthonormal basis in which every diagonal entry of $H$ is zero:
choose a unit vector with expectation zero between its smallest and largest
eigenvalues, restrict to its orthogonal complement, and induct on dimension.
In this basis, writing $\eta=\|H\|>0$, choose
\begin{equation}
 A_{jj}=\sqrt\eta\left(j-\frac{q+1}{2}\right),\qquad
 B_{jh}=\frac{H_{jh}}{i(A_{jj}-A_{hh})}\quad(j\ne h),\qquad B_{jj}=0.
\end{equation}
Then $A,B$ are traceless Hermitian, $i[A,B]=H$, and
$\|A\|,\|B\|\leq C_q\sqrt\eta$. Transporting back to the original
basis preserves these statements. The case $H=0$ uses $A=B=0$.
For $V=e^{iA}$ and $W=e^{iB}$, the exponential series through degree two gives
\begin{equation}\label{eq:finite-gate-balanced-commutator}
 \|VWV^{-1}W^{-1}-\Delta\|
 \leq C_q\|\Delta-I\|^{3/2},\qquad
 \|V-I\|,\|W-I\|\leq C_q\sqrt{\|\Delta-I\|}.
\end{equation}
The remainder bound follows by summing the absolutely convergent series of
total degree at least three. All constants depend only on $q$.

For unitary $V,W,V',W'$ the following elementary estimate retains the
smallness of the commutator factors:
\begin{align}\label{eq:finite-gate-commutator-robustness}
 &\|VWV^{-1}W^{-1}-V'W'V'^{-1}W'^{-1}\|\\
 &\quad\leq2\|V-V'\|\|W-I\|
       +2\|W-W'\|\|V'-I\|.\nonumber
\end{align}
Indeed, varying $V$ is a difference of conjugations of $W-I$, and varying
$W$ is a difference of conjugations of $V'^{-1}-I$. Inverse words are used
exactly, rather than approximating the inverse factors independently.

Density and compactness provide a finite set of words approximating every
special unitary to a fixed sufficiently small error $e_0$; let $L_0$ bound
their lengths. Suppose all special unitaries can be approximated to error
$e_n$ with words of length at most $L_n$. Given $U$, choose such a word $U_n$
and apply \eqref{eq:finite-gate-balanced-commutator} to
$\Delta=UU_n^{-1}$. Approximate its two commutator factors to error $e_n$
using the same previous-level construction, obtaining $V_n,W_n$.
Equations \eqref{eq:finite-gate-balanced-commutator} and
\eqref{eq:finite-gate-commutator-robustness} show that
$V_nW_nV_n^{-1}W_n^{-1}U_n$ approximates $U$ to error
\begin{equation}
 e_{n+1}=C e_n^{3/2},\qquad L_{n+1}=5L_n.
\end{equation}
Choose $e_0<C^{-2}$, decreasing it further for the logarithm construction.
Then $e_n=C^{-2}(C^2e_0)^{(3/2)^n}$ and $L_n=5^nL_0$.
Selecting the first level with $e_n\leq\delta$ proves the asserted bound.
\end{proof}

\begin{proof}[Proof of Theorem~\ref{thm:finite-gate-phase-lift}]
For a local $k$-qubit gate $V\in\operatorname{U}(2^k)$, $k=1,2$, choose
\begin{equation}\label{eq:finite-gate-local-phase-lift}
 W_V=\operatorname{diag}\bigl((\det V)^{-1},1,\ldots,1\bigr),\qquad
 \widehat V=|0\rangle\langle0|\otimes V
             +|1\rangle\langle1|\otimes W_V.
\end{equation}
This is a special unitary on $k+1$ qubits. It sends
$|0\rangle\otimes\psi$ exactly to $|0\rangle\otimes V\psi$, including
the scalar phase of $V$. The same additional qubit is used at every gate.
The ideal lifted product $\widehat C$ therefore satisfies
$\widehat C J_+=J_+U$.

Apply the two lemmas with $\delta=\varepsilon/G$ to each lifted local gate.
Every resulting special-unitary word is an actual alphabet word, with
constant overhead. Unitary telescoping on the complete register gives
\begin{equation}
 \|C_\varepsilon-\widehat C\|
 \leq\sum_{g=1}^G\|\widehat V_{g,\delta}-\widehat V_g\|
 \leq\varepsilon.
\end{equation}
Since $J_+$ is an isometry, this implies \eqref{eq:finite-gate-input-error}.
Applying the orthogonal projection $I-J_+J_+^\dagger$ proves the leakage
bound. The same estimate holds with an arbitrary reference system tensored
with the identity, so it also applies to entangled physical inputs.
No full-register approximation to $U\otimes I$ is claimed: the ideal
$\widehat C$ specifies a different unitary on the unused auxiliary sectors.

For the line construction, interleave one phase qubit after each original
register qubit. A neighboring original pair and the phase qubit between
them form a contiguous three-qubit interval. An original one-qubit gate uses
the adjacent phase qubit. Thus each lifted gate lies in a two- or three-qubit
interval; any nonneighboring CNOT within this interval is replaced by a
constant number of neighboring swaps and CNOTs. All ideal phase qubits stay
zero. The same full-register telescoping argument applies even when earlier
approximation errors have nonzero components in these qubits.
If only one phase qubit is added instead, routing it to each original local
gate and undoing the routing uses $O(m)$ exact swaps per gate.
\end{proof}

\begin{corollary}[Application to bounded-cut-rank MPUs]
An $N$-qubit unitary with consecutive operator Schmidt ranks at most $D$
admits a circuit over \eqref{eq:finite-gate-alphabet} whose initialized-input
operator-norm error is at most $\varepsilon$, retaining the actual phase,
with the arbitrary-gate bound $G$ multiplied by
$O((1+\log(G/\varepsilon))^c)$. With arbitrary pairs, one additional qubit initialized to zero is
sufficient. With neighboring gates, interleaving phase qubits
retains this gate bound and $O(N\log(D+1))$ auxiliary qubits. For fixed $D$
this is polynomial in $N$ and $\log(1/\varepsilon)$.
The approximation allows auxiliary leakage of norm at most $\varepsilon$.
It is a nonuniform existence statement and asserts no efficient classical
computation of the determinant-maximizing metrics or circuit parameters.
\end{corollary}

\section{An obstruction to exact auxiliary erasure}

The allowance for auxiliary error cannot simply be deleted for the stated
alphabet. Let
\begin{equation}
 R=\mathbb Z[1/\sqrt2,i]=\mathbb Z[\sqrt2,i,1/2],\qquad
 \mu_8=\{e^{i\pi h/4}:h\in\mathbb Z\}.
\end{equation}
Every matrix entry of an alphabet circuit belongs to $R$.

\begin{lemma}\label{lem:finite-gate-ring-unit-circle}
The unit-circle elements of $R$ are exactly $\mu_8$.
\end{lemma}

\begin{proof}
Write
\begin{equation}
 z=\frac{a+b\sqrt2+i(c+d\sqrt2)}{2^h},\qquad a,b,c,d\in\mathbb Z,
 \quad h\geq0.
\end{equation}
The identity $|z|=1$ gives
\begin{equation}\label{eq:finite-gate-ring-diophantine}
 a^2+c^2+2b^2+2d^2=2^{2h},\qquad ab+cd=0.
\end{equation}
If $b=d=0$, reduction modulo four and division by four repeatedly show that
$(a,c)=(\pm2^h,0)$ or $(0,\pm2^h)$.
Otherwise write $b=gp,d=gq$ with $g>0$ and $\gcd(p,q)=1$.
The second equation implies $a=tq,c=-tp$ for an integer $t$. Hence
\begin{equation}
 (p^2+q^2)(t^2+2g^2)=2^{2h}.
\end{equation}
The coprimality of $p,q$ shows that $p^2+q^2$ is either $1$ or $2$:
if odd it must be $1$, and if even then $p,q$ are odd and their squared sum
is $2$ modulo eight, so it must be $2$.

For $x^2+2y^2=2^s$ with $s\geq2$, reduction modulo four forces both $x,y$
to be even. Repeated division by four, followed by the cases $s=0,1$,
therefore gives
\begin{align}
 x^2+2y^2=2^{2h}&\ \Longrightarrow\ y=0,\ x=\pm2^h,\\
 x^2+2y^2=2^{2h-1}&\ \Longrightarrow\ x=0,\ y=\pm2^{h-1}
 \quad(h\geq1).
\end{align}
The case $p^2+q^2=1$ contradicts $g>0$.
The other case has $p,q\in\{\pm1\}$, $t=0$, and $g=2^{h-1}$,
giving $z=(\pm1\pm i)/\sqrt2$. Together with the first case these are
exactly $\mu_8$, all of which belong to $R$.
\end{proof}

\newpage
\begin{theorem}[Determinant obstruction with exactly erased auxiliaries]
\label{thm:finite-gate-clean-obstruction}
Let $U\in\operatorname{U}(M)$, where $M=2^N$, and set
$a_U=\operatorname{dist}(\det U,\mu_8)$. If an alphabet circuit $C$ with
any number of clean auxiliary qubits erases them exactly and induces the
physical unitary $V$, then
\begin{equation}\label{eq:finite-gate-clean-lower-bound}
 \|CJ-JU\|=\|V-U\|\geq\frac{a_U}{2^N}.
\end{equation}
In particular, when $\det U\notin\mu_8$, arbitrarily accurate
operator-norm approximation with exact auxiliary erasure is impossible
over this alphabet, regardless of circuit size or auxiliary count.
\end{theorem}

\begin{proof}
Exact erasure says $CJ=JV$. Thus $V=J^\dagger CJ$ has entries in $R$;
it is unitary since $C,J$ are isometries and the physical input and output
spaces have the same dimension. Its determinant is an element of $R$ of
modulus one, hence belongs to $\mu_8$ by
Lemma~\ref{lem:finite-gate-ring-unit-circle}.
For any two $M$-dimensional unitaries $V,U$, let $\lambda_1,\ldots,\lambda_M$
be the eigenvalues of $U^\dagger V$. Then
\begin{equation}
 |\det V-\det U|
 =\left|\prod_{h=1}^M\lambda_h-1\right|
 \leq\sum_{h=1}^M|\lambda_h-1|
 \leq M\|V-U\|.
\end{equation}
This gives the lower bound. Finally $J$ preserves the norm of every vector,
so the initialized-map norm equals $\|V-U\|$.
\end{proof}

The obstruction already occurs at bond dimension one. For example,
$U=e^{i\pi/(8\cdot2^N)}I_{2^N}$ has determinant $e^{i\pi/8}$ and
every operator Schmidt rank equals one. Its approximation error with exact auxiliary erasure over the stated alphabet is at least
$2\sin(\pi/16)/2^N$. This does not contradict
Theorem~\ref{thm:finite-gate-phase-lift}, which allows a residual auxiliary
component. The unused sector of the lifted unitary carries the compensating
determinant, and finite-word approximation need not preserve its invariant
subspaces exactly.

\section{Exact-erasure approximation in a fixed physical dimension}

Fix $N\geq1$ and $M=2^N$. Let $\mathcal E_N$ be the physical unitaries
$V\in\operatorname{U}(M)$ admitting an actual alphabet circuit $C$ on
$N+a$ qubits, for some finite $a\geq0$, with $CJ_a=J_aV$ exactly, where
$J_a\psi=\psi\otimes|0\rangle^{\otimes a}$. Thus every auxiliary is
returned to zero on every physical input. The following closure is taken
in phase-sensitive operator norm; the auxiliary count may vary in
$\mathcal E_N$.

\begin{theorem}[Closure of exactly erased implementations]
\label{thm:finite-gate-clean-closure}
For every fixed $N\geq1$,
\[
 \overline{\mathcal E_N}
 =\{U\in\operatorname{U}(2^N):\det U\in\mu_8\}.
\]
If $U\in\operatorname{SU}(2^N)$, approximating circuits may act on the
physical qubits alone. If $\det U\in\mu_8$, at most $N-1$ clean auxiliary
qubits suffice. For $0<\varepsilon<1$, the gate count is at most
\[
 C_N\bigl(1+\log(1/\varepsilon)\bigr)^c,
 \qquad c=\frac{\log5}{\log(3/2)},
\]
where $C_N$ depends on $N$. No uniform polynomial bound in $N$ is asserted
for these constants.
\end{theorem}

\begin{proof}
Necessity follows from Lemma~\ref{lem:finite-gate-ring-unit-circle}
and continuity: every $V\in\mathcal E_N$ has entries in $R$ and determinant
in $\mu_8$.

For sufficiency, extend Lemma~\ref{lem:finite-gate-su-density} to $N$
physical qubits. Take its local words $P_i,Q_i$ and, for $i\ne j$,
\[
 \operatorname{CNOT}_{ij}P_j\operatorname{CNOT}_{ij},\qquad
 \operatorname{CNOT}_{ij}Q_j\operatorname{CNOT}_{ij},
\]
together with their inverses; for $N=1$, use only $P,Q$ and their inverses.
This inverse-closed set $\mathcal S_N$ consists of exact constant-length
alphabet words in $\operatorname{SU}(2^N)$. Its closed generated group
contains every local special unitary and, by
$\operatorname{CNOT}_{ij}Z_j\operatorname{CNOT}_{ij}=Z_iZ_j$, all rotations
generated by $Z_iZ_j$. Local conjugations give all two-site Pauli generators.

Inductively, $i$ times every nonidentity Pauli product lies in the generated
Lie algebra. To extend a product to one further site, choose two distinct
Pauli factors at a common site whose product is $\pm i$ times the desired
factor there. The first generator has the existing support and the second
has the common site and the new site. Their commutator is a nonzero real
multiple of $i$ times the desired larger product. The resulting Pauli
products span the traceless Hermitian matrices. Product and commutator
limits give their exponentials; every special unitary has a traceless
logarithm by choosing eigenangles with sum zero. Thus the subgroup generated
by $\mathcal S_N$ is dense in $\operatorname{SU}(2^N)$ with the actual phase
and no auxiliaries.
Lemma~\ref{lem:finite-gate-quantitative-su}, with $q=2^N$, gives the asserted
word bound. Its initial net and commutator constants depend on $N$.

It remains to implement each determinant coset. For $h\in\mathbb Z$, put
\[
 D_h=\operatorname{diag}(\zeta^h,1,\ldots,1),\qquad \zeta=e^{i\pi/4}.
\]
An exact $O(N)$-gate circuit uses at most $N-1$ clean auxiliaries. Apply
$X=HT^4H$ to each physical qubit and, for $N\geq2$, compute successive
conjunctions into $N-1$ zero bits with Toffoli gates. The last bit is one
exactly when the original configuration was all zero. Apply $T^h$ there,
then undo the conjunctions and the $X$ gates. The phase gate changes no
computational value, so all auxiliaries return to zero. For $N=1$, use
$XT^hX$. Powers of $T$ have bounded length by reduction modulo eight.

Here Toffoli is an exact alphabet word. The integer identity
\[
 4xyz=x+y+z-(x\mathbin\oplus y)-(x\mathbin\oplus z)
       -(y\mathbin\oplus z)+(x\mathbin\oplus y\mathbin\oplus z)
 \quad(x,y,z\in\{0,1\})
\]
implements controlled-controlled-$Z$: compute each parity into one of the
three qubits by CNOTs, apply $T$ or $T^{-1}$, and uncompute. The phase is
exactly $\zeta^{4xyz}=(-1)^{xyz}$. Conjugating the target by $H$ gives
Toffoli, with no auxiliary or extra scalar phase. The conjunction circuit
therefore erases exactly on coherent superpositions.

If $\det U=\zeta^h$, approximate $D_h^\dagger U\in\operatorname{SU}(M)$
by a physical $\mathcal S_N$-word, followed by the exact clean $D_h$
circuit. The error is unchanged and every auxiliary is erased. The bounded
length of each generator and the $O(N)$ coset circuit are absorbed in $C_N$.
\end{proof}

\begin{corollary}[Sharp distance to exactly erased implementations]
\label{cor:finite-gate-clean-sharp-distance}
Let $\delta_U\in[0,\pi/8]$ be the least angular distance from $\det U$ to
$\mu_8$, where $U\in\operatorname{U}(M)$ and $M=2^N$. Then
\[
 \inf_{V\in\mathcal E_N}\|V-U\|
 =2\sin\left(\frac{\delta_U}{2M}\right).
\]
\end{corollary}

\begin{proof}
For $t=\|V-U\|$, the eigenangles $\phi_j\in[-\pi,\pi]$ of $U^\dagger V$
satisfy $|\phi_j|\leq2\arcsin(t/2)$. Their sum represents the relative
determinant angle modulo $2\pi$, so
$\delta_U\leq2M\arcsin(t/2)$. This proves the lower bound.
Choose $\zeta^h$ with $\det U/\zeta^h=e^{i\delta}$ and
$|\delta|=\delta_U$. The unitary $V_0=e^{-i\delta/M}U$ has determinant
$\zeta^h$ and distance $2\sin(\delta_U/(2M))$ from $U$.
Theorem~\ref{thm:finite-gate-clean-closure} approximates $V_0$ arbitrarily
closely by members of $\mathcal E_N$, proving the upper bound for the infimum.
\end{proof}

The scalar example above therefore has exact best error
$2\sin(\pi/(16\cdot2^N))$, attained by the identity circuit. Its earlier
lower bound remains valid but is weaker. For special-unitary targets,
exact-erasure approximation is possible even without auxiliaries; the
constants above do not establish a polynomial-in-$N$ gate bound for
bounded-cut-rank targets. Allowing small auxiliary leakage retains the arbitrary-gate MPU bound,
multiplied by the precision factor in Section~1.

\begin{thebibliography}{1}
\bibitem{exact-mpu-circuits}
TNLean contributors,
\emph{Exact polynomial circuits for matrix product unitaries},
2 October 2026, Section~5, \texttt{thm:general-determinant-circuits}.
Local source:
\nolinkurl{docs/audits/2026-10-02_mpu_rank_two_circuits.tex}.
\end{thebibliography}
\end{document}
